% ======================================================================
\chapter[Радиоактивность. Закон распада. Вековое уравнение]%
        {Ядерная радиоактивность. Закон радиоактивного распада. Вековое уравнение}
\label{t:18}
% ======================================================================

\begin{qbox}
\textbf{Билет 18.} Ядерная радиоактивность. Закон радиоактивного распада, константа распада,
период полураспада, среднее время жизни, вековое уравнение.
\end{qbox}

\section{Содержание билета}

\textbf{Подтемы:}
\begin{enumerate}
 \item явление радиоактивности, виды распадов, их общие свойства;
 \item закон радиоактивного распада: вывод из статистической гипотезы, постоянная распада;
 \item период полураспада, среднее время жизни и связь между ними; активность и её единицы;
 \item последовательные распады (цепочки): решение уравнений, вековое (секулярное) и подвижное равновесие; вековое уравнение;
 \item применения: радиоактивные семейства, датирование, определение больших периодов полураспада.
\end{enumerate}

\textbf{Необходимые понятия:} энергия связи и стабильность ядер (\tref{t:17}); вероятность, распределение
Пуассона; линейные дифференциальные уравнения первого порядка; ширина уровня и время жизни (\tref{t:02}).

\section{Определения и физические принципы}

\begin{definition}[радиоактивность]
Самопроизвольное (спонтанное) превращение неустойчивых атомных ядер в другие ядра, сопровождающееся испусканием
частиц или излучения. Основные виды: $\alpha$-распад (испускание ядра $\nucl{He}{4}$), $\beta$-распад ($\beta^-$, $\beta^+$, электронный
захват), $\gamma$-излучение (переход ядра между уровнями), спонтанное деление; реже — испускание протонов, кластеров.
\end{definition}

\begin{definition}[постоянная распада]
\emph{Постоянная распада} $\lambda$ — вероятность распада одного ядра в единицу времени: вероятность распасться за малый
интервал $\dd t$ равна $\lambda\,\dd t$. Размерность $[\lambda]=\text{с}^{-1}$.
\end{definition}

\begin{definition}[период полураспада, среднее время жизни, активность]
\emph{Период полураспада} $T_{1/2}$ — время, за которое распадается половина имеющихся ядер. \emph{Среднее время жизни}
$\tau$ — среднее значение времени существования ядра до распада. \emph{Активность} $\mathcal A$ — число распадов в единицу
времени: $\mathcal A=-\dd N/\dd t=\lambda N$. Единицы: беккерель (1~Бк $=1$ распад/с) и кюри (1~Ки $=3{,}7\cdot10^{10}$~Бк —
примерно активность 1~г $\nucl{Ra}{226}$).
\end{definition}

\emph{Физический смысл.} Радиоактивный распад — квантовый переход из квазистационарного состояния; момент распада
отдельного ядра принципиально непредсказуем, определена лишь вероятность. Постоянная $\lambda$ — свойство ядра:
она не зависит от времени существования ядра (ядра «не стареют»), от температуры, давления, химического состояния
(за исключением слабой зависимости для электронного захвата и внутренней конверсии, где важна плотность электронов
у ядра, — эффекты порядка $10^{-3}$).

\emph{Область применимости закона.} Экспоненциальный закон описывает поведение большого числа ядер в среднем; при
малых $N$ существенны флуктуации (распределение Пуассона, относительная погрешность $1/\sqrt{N_{\text{расп}}}$).

\section{Основная теория}

\subsection{Открытие и виды распадов}

А.~Беккерель (1896) обнаружил излучение солей урана; М.~и П.~Кюри открыли полоний и радий; Э.~Резерфорд разделил
излучение на $\alpha$, $\beta$, $\gamma$ по отклонению в магнитном поле и установил (вместе с Ф.~Содди, 1902), что радиоактивность —
превращение одних элементов в другие. Правила смещения: при $\alpha$-распаде $(A,Z)\to(A-4,Z-2)$, при $\beta^-$-распаде
$(A,Z)\to(A,Z+1)$, при $\beta^+$-распаде и электронном захвате $(A,Z)\to(A,Z-1)$, при $\gamma$-переходе $A$ и $Z$ не меняются.

Периоды полураспада лежат в огромном диапазоне: от $\sim10^{-22}$~с до $\gtrsim10^{24}$~лет ($\nucl{Te}{128}$, двойной бета-распад).
Причины: экспоненциальная зависимость вероятности туннелирования для $\alpha$-распада (\tref{t:19}), сильная зависимость
от энергии ($\propto Q^5$) и правил отбора для $\beta$-распада (\tref{t:20}).

\subsection{Закон радиоактивного распада}

Из определения постоянной распада (вывод~\ref{der:18:law}):
\begin{equation}\label{eq:18:law}
  \diff{N}{t}=-\lambda N\quad\Longrightarrow\quad N(t)=N_0\,\ee^{-\lambda t}=N_0\,2^{-t/T_{1/2}} .
\end{equation}
Связь характерных времён (вывод~\ref{der:18:tau}):
\begin{equation}\label{eq:18:times}
  T_{1/2}=\frac{\ln2}{\lambda}\approx\frac{0{,}693}{\lambda},\qquad \tau=\frac1\lambda,\qquad T_{1/2}=\tau\ln2 .
\end{equation}
За время $\tau$ число ядер уменьшается в $\ee$ раз. Активность убывает по тому же закону: $\mathcal A(t)=\lambda N_0\ee^{-\lambda t}$.
Измеряя $\mathcal A$ в разные моменты, определяют $\lambda$ по наклону прямой $\ln\mathcal A(t)$ (для коротких периодов); для
очень больших $T_{1/2}$ используют $\lambda=\mathcal A/N$ при известном числе ядер.

\paragraph{Связь с шириной уровня.} Квазистационарное состояние со временем жизни $\tau$ имеет энергетическую ширину
$\Gamma=\hbar/\tau=\hbar\lambda$; $|\Psi|^2\propto\ee^{-t/\tau}$. Это связывает закон распада с соотношением энергия--время (\tref{t:02}).

\paragraph{Несколько каналов распада.} Если ядро распадается несколькими способами с постоянными $\lambda_i$, то полная
постоянная $\lambda=\sum_i\lambda_i$ (вероятности несовместных событий складываются), а доля канала $i$ (коэффициент
ветвления) равна $\lambda_i/\lambda$. Время жизни определяется полной $\lambda$.

\begin{figure}[htbp]
\centering
\begin{tikzpicture}
\begin{axis}[width=0.66\linewidth, height=5.0cm, xlabel={$t/T_{1/2}$}, ylabel={$N/N_0$}, xmin=0, xmax=4.2, ymin=0, ymax=1.08,
   samples=150, xtick={0,1,1.4427,2,3,4}, xticklabels={$0$,$1$,$\tau$,$2$,$3$,$4$}, ytick={0,0.125,0.25,0.3679,0.5,1},
   yticklabels={$0$,$\frac18$,$\frac14$,$\frac1{\ee}$,$\frac12$,$1$}, grid=major, grid style={gray!25}]
  \addplot[very thick, accent, domain=0:4.2] {2^(-x)};
\end{axis}
\end{tikzpicture}
\caption{Закон радиоактивного распада $N/N_0=2^{-t/T_{1/2}}$. Отмечены период полураспада и среднее время жизни
$\tau=T_{1/2}/\ln2\approx1{,}44\,T_{1/2}$.}
\label{fig:18:decay}
\end{figure}

\subsection{Последовательные распады}

Пусть материнское ядро 1 распадается в дочернее 2, которое тоже радиоактивно: $1\xrightarrow{\lambda_1}2\xrightarrow{\lambda_2}3$.
Система уравнений
\begin{equation}\label{eq:18:chain}
  \diff{N_1}{t}=-\lambda_1N_1,\qquad \diff{N_2}{t}=\lambda_1N_1-\lambda_2N_2 .
\end{equation}
Первое слагаемое во втором уравнении — образование дочерних ядер, второе — их распад. При $N_2(0)=0$
(вывод~\ref{der:18:chain})
\begin{equation}\label{eq:18:N2}
  N_2(t)=N_{10}\frac{\lambda_1}{\lambda_2-\lambda_1}\Bigl(\ee^{-\lambda_1t}-\ee^{-\lambda_2t}\Bigr).
\end{equation}
Число дочерних ядер сначала растёт, достигает максимума при $t_m=\dfrac{\ln(\lambda_2/\lambda_1)}{\lambda_2-\lambda_1}$ (в этот момент
$\lambda_1N_1=\lambda_2N_2$ — активности равны), а затем поведение зависит от соотношения постоянных.

\paragraph{1. Вековое (секулярное) равновесие: $\lambda_1\ll\lambda_2$} ($T_1\gg T_2$). При $t\gg\tau_2$ в~\eqref{eq:18:N2} $\ee^{-\lambda_2t}\to0$,
$\lambda_2-\lambda_1\approx\lambda_2$, и $N_2\approx(\lambda_1/\lambda_2)N_1$. Получаем \emph{вековое уравнение}
\begin{equation}\label{eq:18:secular}
  \boxed{\lambda_1N_1=\lambda_2N_2}\qquad\Longleftrightarrow\qquad \frac{N_1}{T_1}=\frac{N_2}{T_2},\qquad \mathcal A_1=\mathcal A_2 .
\end{equation}
Активности материнского и дочернего нуклидов равны: дочерних ядер образуется столько же, сколько распадается.
Равновесие устанавливается за несколько $T_2$ и сохраняется сколь угодно долго (на масштабах $\ll T_1$) — отсюда
название «вековое».

\paragraph{2. Подвижное (переходное) равновесие: $\lambda_1<\lambda_2$, но сравнимы.} При $t\gg1/(\lambda_2-\lambda_1)$
$N_2/N_1\to\lambda_1/(\lambda_2-\lambda_1)$ — постоянное отношение; активность дочернего нуклида больше материнской:
$\mathcal A_2/\mathcal A_1=\lambda_2/(\lambda_2-\lambda_1)>1$, и оба спадают с периодом материнского.

\paragraph{3. Нет равновесия: $\lambda_1>\lambda_2$.} Материнский нуклид быстро распадается, затем остаются только дочерние
ядра, распадающиеся со своим периодом.

\paragraph{Цепочка из многих звеньев.} Для $1\to2\to\dots\to k$ при $\lambda_1\ll\lambda_i$ ($i\ge2$) через время, большое по сравнению
со всеми $T_i$ ($i\ge2$), устанавливается равновесие всего семейства:
\begin{equation}\label{eq:18:family}
  \lambda_1N_1=\lambda_2N_2=\dots=\lambda_kN_k .
\end{equation}

\begin{figure}[htbp]
\centering
\begin{tikzpicture}
\begin{axis}[width=0.48\linewidth, height=5.0cm, xlabel={$t$ (отн. ед.)}, ylabel={активность}, xmin=0, xmax=10, ymin=0, ymax=1.12,
   samples=200, title={\small вековое: $\lambda_1=0{,}01$, $\lambda_2=1$}, legend style={font=\scriptsize, draw=none}, legend pos=south east]
  \addplot[very thick, accent, domain=0:10] {exp(-0.01*x)};
  \addlegendentry{$\lambda_1N_1$}
  \addplot[very thick, keycol, dashed, domain=0:10] {1/(1-0.01)*(exp(-0.01*x)-exp(-x))};
  \addlegendentry{$\lambda_2N_2$}
\end{axis}
\end{tikzpicture}\hfill
\begin{tikzpicture}
\begin{axis}[width=0.48\linewidth, height=5.0cm, xlabel={$t$ (отн. ед.)}, ylabel={активность}, xmin=0, xmax=10, ymin=0, ymax=1.12,
   samples=200, title={\small подвижное: $\lambda_1=0{,}3$, $\lambda_2=1$}, legend style={font=\scriptsize, draw=none}, legend pos=north east]
  \addplot[very thick, accent, domain=0:10] {exp(-0.3*x)};
  \addlegendentry{$\lambda_1N_1$}
  \addplot[very thick, keycol, dashed, domain=0:10] {1/(1-0.3)*(exp(-0.3*x)-exp(-x))};
  \addlegendentry{$\lambda_2N_2$}
\end{axis}
\end{tikzpicture}
\caption{Активности материнского и дочернего нуклидов (расчёт по~\eqref{eq:18:N2}, $N_2(0)=0$, активности в единицах
$\lambda_1N_{10}$). Слева — установление векового равновесия $\lambda_1N_1=\lambda_2N_2$; справа — подвижное равновесие,
при котором активность дочернего нуклида превышает материнскую.}
\label{fig:18:chain}
\end{figure}

\subsection{Применения}

\paragraph{Радиоактивные семейства.} Тяжёлые природные нуклиды образуют четыре семейства, различающиеся остатком
от деления $A$ на 4 (при $\alpha$-распаде $A$ меняется на 4, при $\beta$-распаде не меняется): $4n$ (родоначальник $\nucl{Th}{232}$,
$T=1{,}4\cdot10^{10}$~лет), $4n+2$ ($\nucl{U}{238}$, $4{,}47\cdot10^9$~лет), $4n+3$ ($\nucl{U}{235}$, $7{,}0\cdot10^8$~лет), $4n+1$ ($\nucl{Np}{237}$,
$2{,}1\cdot10^6$~лет — в природе практически вымерло). Ряд урана заканчивается стабильным $\nucl{Pb}{206}$ после 8 $\alpha$- и
6 $\beta^-$-распадов. В урановой руде все члены ряда находятся в вековом равновесии.

\paragraph{Определение больших периодов полураспада.} Измерив отношение количеств двух членов ряда в равновесной руде и зная
период одного из них, находят период другого по~\eqref{eq:18:secular} (вывод~\ref{der:18:RaU}).

\paragraph{Датирование.} Радиоуглеродный метод: в живых организмах отношение $\nucl{C}{14}/\nucl{C}{12}$ поддерживается постоянным
за счёт обмена с атмосферой (где $\nucl{C}{14}$ образуется космическими лучами); после гибели $\nucl{C}{14}$ распадается с
$T_{1/2}=5730$~лет, и возраст
\begin{equation}\label{eq:18:dating}
  t=\frac{T_{1/2}}{\ln2}\ln\frac{\mathcal A_0}{\mathcal A}.
\end{equation}
Геологический возраст определяют по отношениям $\nucl{Pb}{206}/\nucl{U}{238}$, $\nucl{Ar}{40}/\nucl{K}{40}$ и др.

\section{Математические выводы}

\begin{derivation}[label=der:18:law]{закон радиоактивного распада}
\emph{Допущения:} (1) ядра распадаются независимо друг от друга; (2) вероятность распада ядра за время $\dd t$ равна $\lambda\dd t$ и не
зависит от того, сколько времени ядро уже существует; (3) $N$ велико.

За время $\dd t$ из $N$ ядер в среднем распадается $\dd N_{\text{расп}}=N\lambda\dd t$, поэтому изменение числа ядер
\[
  \dd N=-\lambda N\dd t\quad\Longrightarrow\quad\frac{\dd N}{N}=-\lambda\dd t\quad\Longrightarrow\quad\ln N=-\lambda t+\const .
\]
С начальным условием $N(0)=N_0$: $N(t)=N_0\ee^{-\lambda t}$.

\emph{Вероятностная форма.} Вероятность того, что отдельное ядро не распадётся до момента $t$: $P(t)$. По допущению (2)
$P(t+\dd t)=P(t)(1-\lambda\dd t)$, откуда $\dd P/\dd t=-\lambda P$ и $P(t)=\ee^{-\lambda t}$. Экспоненциальный закон — единственное
распределение времени жизни «без памяти»: $P(t+s\,|\,\text{дожило до }s)=P(t)$.
\end{derivation}

\begin{derivation}[label=der:18:tau]{период полураспада и среднее время жизни}
\emph{Период полураспада.} $N(T_{1/2})=N_0/2$: $\ee^{-\lambda T_{1/2}}=\frac12$ $\Rightarrow$ $T_{1/2}=\ln2/\lambda$.

\emph{Среднее время жизни.} Плотность вероятности распада в момент $t$: $w(t)=-\dd P/\dd t=\lambda\ee^{-\lambda t}$ (нормирована:
$\int_0^\infty w\dd t=1$). Среднее
\[
  \tau=\int_0^\infty t\,\lambda\ee^{-\lambda t}\dd t=\Bigl[-t\ee^{-\lambda t}\Bigr]_0^\infty+\int_0^\infty\ee^{-\lambda t}\dd t=0+\frac1\lambda=\frac1\lambda .
\]
Следовательно, $T_{1/2}=\tau\ln2\approx0{,}693\,\tau$. Эквивалентно: $\tau=\frac{1}{N_0}\int_0^\infty N(t)\dd t$ — «суммарное время жизни всех ядер,
делённое на их число».
\end{derivation}

\begin{derivation}[label=der:18:chain]{решение уравнений цепочки $1\to2\to3$}
\emph{Шаг 1.} $N_1=N_{10}\ee^{-\lambda_1t}$. Подставим во второе уравнение~\eqref{eq:18:chain}:
$\dot N_2+\lambda_2N_2=\lambda_1N_{10}\ee^{-\lambda_1t}$.

\emph{Шаг 2 (интегрирующий множитель $\ee^{\lambda_2t}$).}
\[
  \diff{}{t}\bigl(N_2\ee^{\lambda_2t}\bigr)=\lambda_1N_{10}\ee^{(\lambda_2-\lambda_1)t}
  \ \Rightarrow\
  N_2\ee^{\lambda_2t}=\frac{\lambda_1N_{10}}{\lambda_2-\lambda_1}\ee^{(\lambda_2-\lambda_1)t}+C .
\]
\emph{Шаг 3 (начальное условие $N_2(0)=0$):} $C=-\dfrac{\lambda_1N_{10}}{\lambda_2-\lambda_1}$, и
\[
  N_2(t)=\frac{\lambda_1N_{10}}{\lambda_2-\lambda_1}\Bigl(\ee^{-\lambda_1t}-\ee^{-\lambda_2t}\Bigr).
\]
\emph{Шаг 4 (максимум).} $\dot N_2=0$: $\lambda_1\ee^{-\lambda_1t}=\lambda_2\ee^{-\lambda_2t}$ $\Rightarrow$ $t_m=\dfrac{\ln(\lambda_2/\lambda_1)}{\lambda_2-\lambda_1}$.
В этот момент из~\eqref{eq:18:chain} $\lambda_1N_1=\lambda_2N_2$.

\emph{Шаг 5 (асимптотики).} Отношение активностей
\[
  \frac{\lambda_2N_2}{\lambda_1N_1}=\frac{\lambda_2}{\lambda_2-\lambda_1}\Bigl(1-\ee^{-(\lambda_2-\lambda_1)t}\Bigr)
  \xrightarrow[t\to\infty]{}\frac{\lambda_2}{\lambda_2-\lambda_1}\quad(\lambda_2>\lambda_1).
\]
При $\lambda_1\ll\lambda_2$ предел равен 1 — вековое равновесие~\eqref{eq:18:secular}. Предел $\lambda_1\to\lambda_2$ даёт
$N_2=\lambda_1N_{10}t\ee^{-\lambda_1t}$ (раскрытие неопределённости).

\emph{Вековое уравнение напрямую:} если $N_1$ практически постоянно на временах $\sim\tau_2$, то стационарное решение
второго уравнения~\eqref{eq:18:chain} ($\dot N_2=0$) — $\lambda_1N_1=\lambda_2N_2$; к нему решение приближается как $1-\ee^{-\lambda_2t}$.
\end{derivation}

\begin{derivation}[label=der:18:RaU]{отношение радия к урану в руде}
В равновесном ряду урана $\lambda_{\mathrm U}N_{\mathrm U}=\lambda_{\mathrm{Ra}}N_{\mathrm{Ra}}$, т.\,е.
\[
  \frac{N_{\mathrm{Ra}}}{N_{\mathrm U}}=\frac{T_{\mathrm{Ra}}}{T_{\mathrm U}}=\frac{1600\ \text{лет}}{4{,}47\cdot10^9\ \text{лет}}=3{,}6\cdot10^{-7}.
\]
В 1~т урана содержится $\approx10^6\cdot3{,}6\cdot10^{-7}\cdot\frac{226}{238}\approx0{,}34$~г радия. Наоборот, измерив отношение, можно найти неизвестный
период полураспада урана — время, недоступное прямому наблюдению за убыванием активности.

\emph{Активность 1~г радия:} $\lambda=\frac{0{,}693}{1600\cdot3{,}16\cdot10^7\ \text{с}}=1{,}37\cdot10^{-11}$~с$^{-1}$,
$N=\frac{6{,}02\cdot10^{23}}{226}=2{,}66\cdot10^{21}$, $\mathcal A=\lambda N=3{,}6\cdot10^{10}$~Бк $\approx1$~Ки.
\end{derivation}

\section{Итоговые формулы}

\keyformula{N(t)=N_0\ee^{-\lambda t},\qquad \mathcal A=-\diff Nt=\lambda N}
 {$\lambda$ — постоянная распада, $N_0$ — начальное число ядер, $\mathcal A$ — активность.}
 {большое число независимых ядер; $\lambda$ не зависит от времени.}
 {число ядер и активность в любой момент.}

\keyformula{T_{1/2}=\frac{\ln2}{\lambda},\qquad \tau=\frac1\lambda=\frac{T_{1/2}}{\ln2},\qquad \Gamma=\frac{\hbar}{\tau}}
 {$T_{1/2}$ — период полураспада, $\tau$ — среднее время жизни, $\Gamma$ — ширина уровня.}
 {экспоненциальный распад.}
 {пересчёт характерных времён; связь со спектральной шириной.}

\keyformula{N_2=N_{10}\frac{\lambda_1}{\lambda_2-\lambda_1}\bigl(\ee^{-\lambda_1t}-\ee^{-\lambda_2t}\bigr),\qquad \lambda_1N_1=\lambda_2N_2\ (\lambda_1\ll\lambda_2,\ t\gg\tau_2)}
 {индексы 1 и 2 — материнский и дочерний нуклиды; $N_2(0)=0$.}
 {вековое уравнение — при $T_1\gg T_2$ и времени наблюдения $\gg T_2$.}
 {накопление дочерних продуктов; равновесные отношения в рудах; определение больших $T_{1/2}$.}

\keyformula{t=\frac{T_{1/2}}{\ln2}\ln\frac{\mathcal A_0}{\mathcal A},\qquad \lambda=\sum_i\lambda_i}
 {$\mathcal A_0$ — начальная удельная активность, $\lambda_i$ — парциальные постоянные.}
 {известна начальная активность (радиоуглерод: равновесие с атмосферой).}
 {датирование; ветвление распадов.}

\section{Физическая интерпретация}

\textbf{Статистическая природа.} Экспоненциальный закон — следствие независимости распадов и отсутствия «памяти»;
он не означает, что каждое ядро живёт время $\tau$: время жизни отдельного ядра случайно и распределено
экспоненциально. Отклонения от экспоненты предсказываются квантовой механикой на очень малых и очень больших
временах, но на практике ничтожны.

\textbf{Равновесие как баланс потоков.} Вековое уравнение — аналог стационарного уровня жидкости в сосуде с притоком и
стоком: количество короткоживущего продукта устанавливается таким, чтобы скорость его распада равнялась скорости
образования. Поэтому короткоживущие нуклиды (радий, радон) присутствуют в природе, хотя их периоды много меньше
возраста Земли.

\textbf{Масштабы.} $T_{1/2}$ различаются на десятки порядков; при этом постоянная $\lambda$ — одна из наиболее стабильных величин:
она не меняется при изменении внешних условий в химически и физически доступных пределах, что делает распад
надёжными «часами».

\section{Что необходимо знать на экзамене}

\begin{examcheck}
\checkhead{Определения}
\begin{itemize}
 \item радиоактивность, виды распадов, правила смещения; постоянная распада, $T_{1/2}$, $\tau$, активность (Бк, Ки).
\end{itemize}
\checkhead{Формулы}
\begin{itemize}
 \item $N=N_0\ee^{-\lambda t}$; $T_{1/2}=\ln2/\lambda$; $\tau=1/\lambda$; $\mathcal A=\lambda N$; $\lambda=\sum\lambda_i$;
 \item решение для цепочки $N_2(t)$; вековое уравнение $\lambda_1N_1=\lambda_2N_2$; подвижное равновесие.
\end{itemize}
\checkhead{Выводы}
\begin{itemize}
 \item закон распада из гипотезы постоянной вероятности; $\tau=1/\lambda$ через интеграл;
 \item решение системы для $1\to2\to3$ интегрирующим множителем; три режима; $t_m$.
\end{itemize}
\checkhead{Интерпретации}
\begin{itemize}
 \item почему распад случаен и почему $\lambda$ не зависит от внешних условий;
 \item радиоактивные семейства ($4n$, $4n+1$, $4n+2$, $4n+3$); отношение Ra/U; радиоуглеродное датирование.
\end{itemize}
\end{examcheck}
